\documentclass[11pt,letterpaper]{article}
% Reusable style. Build each .tex from its own directory using scripts/build.sh.
\usepackage[T1]{fontenc}
\usepackage[utf8]{inputenc}
\usepackage{lmodern}
\usepackage[margin=0.86in,headheight=15pt]{geometry}
\usepackage{amsmath,amssymb,amsthm,mathtools}
\usepackage{booktabs,longtable,tabularx,array,multicol}
\usepackage[table]{xcolor}
\usepackage{microtype}
\usepackage{enumitem}
\usepackage{fancyhdr}
\usepackage{titlesec}
\usepackage{hyperref}
\usepackage{listings}
\usepackage[most]{tcolorbox}
\usepackage{needspace}
\definecolor{aeNavy}{HTML}{172D43}
\definecolor{aeBlue}{HTML}{176B90}
\definecolor{aeTeal}{HTML}{148477}
\definecolor{aePale}{HTML}{F2F7F9}
\definecolor{aeGray}{HTML}{5C6771}
\hypersetup{colorlinks=true,linkcolor=aeBlue,urlcolor=aeBlue,pdfauthor={Algorithm Engineering Documentation Suite},pdfsubject={Independent educational materials based on the supplied TAOCP volume topics}}
\titleformat{\section}{\Large\bfseries\color{aeNavy}}{\thesection}{0.6em}{}
\titleformat{\subsection}{\large\bfseries\color{aeBlue}}{\thesubsection}{0.6em}{}
\titleformat{\subsubsection}{\normalsize\bfseries\color{aeNavy}}{\thesubsubsection}{0.6em}{}
\setlist[itemize]{leftmargin=*,topsep=3pt,itemsep=1pt}
\setlist[enumerate]{leftmargin=*,topsep=3pt,itemsep=3pt}
\setlength{\parindent}{0pt}\setlength{\parskip}{5pt}
\pagestyle{fancy}\fancyhf{}\fancyhead[L]{\footnotesize\color{aeGray}ALGORITHM ENGINEERING}\fancyhead[R]{\footnotesize\color{aeGray}DOCUMENTATION SUITE}\fancyfoot[L]{\footnotesize\color{aeGray}Independent study and engineering practice}\fancyfoot[R]{\footnotesize\thepage}
\renewcommand{\headrulewidth}{0.35pt}
\newtcolorbox{engineeringnote}[1][]{breakable,colback=aePale,colframe=aeBlue,boxrule=0.5pt,arc=1mm,fonttitle=\bfseries,title={Engineering note},#1}
\newtcolorbox{evidencebox}[1][]{breakable,colback=white,colframe=aeTeal,boxrule=0.6pt,arc=1mm,fonttitle=\bfseries,title={Evidence and deliverables},#1}
\lstdefinestyle{algo}{basicstyle=\ttfamily\footnotesize,breaklines=true,frame=single,rulecolor=\color{aeGray!40},backgroundcolor=\color{aePale},showstringspaces=false,tabsize=2,columns=fullflexible,keepspaces=true,keywordstyle=\color{aeBlue}\bfseries,commentstyle=\color{aeTeal}}
\lstset{style=algo}
\newcommand{\DocTitle}[3]{%
\hypersetup{pageanchor=false}\begin{titlepage}\vspace*{1.2cm}
{\color{aeTeal}\rule{\linewidth}{2.3pt}}\par\vspace{1.15cm}
{\fontsize{24}{30}\selectfont\bfseries\color{aeNavy}\raggedright #1\par}\vspace{0.45cm}
{\Large\color{aeBlue}#2\par}\vspace{1.3cm}
{\large\bfseries #3\par}\vspace{0.35cm}
{\color{aeGray}Edition 1.0 \quad | \quad 29 October 2026\par}\vfill
{\color{aeGray}\small Original educational material aligned to the chapter and section topics in the five supplied volumes of Donald E. Knuth's \emph{The Art of Computer Programming}. This document is not an authorized reproduction or substitute for the original books.}\par\vspace{0.4cm}
{\color{aeTeal}\rule{\linewidth}{0.75pt}}
\end{titlepage}\hypersetup{pageanchor=true}\pagenumbering{roman}}
\newcommand{\SourceNote}[1]{\par\smallskip{\footnotesize\color{aeGray}\textbf{Source locator:} #1}\par\smallskip}
\newcommand{\Outcome}[1]{\par\textbf{Competency target.} #1\par}
\newcommand{\Checkpoint}[1]{\par\textbf{Checkpoint.} #1\par}
% Intentionally simple printable writing lines: call \AnswerLines{number}.
\newcount\alcount
\newcommand{\AnswerLines}[1]{\par\alcount=0\loop\ifnum\alcount<#1\par\noindent\rule{\linewidth}{0.3pt}\vspace{0.33cm}\advance\alcount by1\repeat\par}
\newcommand{\FormField}[2][0.62\linewidth]{\noindent\textbf{#2}\hspace{0.65em}\rule{#1}{0.4pt}\par\vspace{0.18cm}}
\newcommand{\ReviewTable}{\begin{tabularx}{\linewidth}{@{}p{0.23\linewidth}X p{0.2\linewidth}@{}}\toprule Criterion & Recorded evidence & Rating (0--4)\\\midrule Correctness & & \\[1em] Proof quality & & \\[1em] Complexity & & \\[1em] Testing & & \\[1em] Reproducibility & & \\[1em] Communication & & \\[1em]\bottomrule\end{tabularx}}

\begin{document}
\DocTitle{Proof and Complexity Field Guide}{A worked framework for rigorous analysis}{Engineering documentation}
\tableofcontents
\clearpage
\pagenumbering{arabic}
\section{Proof obligation inventory}
\begin{enumerate}\item \textbf{Defined domain:} valid types, sizes, order and numerical ranges.
\item \textbf{Partial correctness:} every return satisfies the postcondition if the algorithm terminates.
\item \textbf{Termination:} a well-founded measure decreases on each progress step.
\item \textbf{Total correctness:} partial correctness together with termination.
\item \textbf{Resource use:} cost model, upper/lower bounds and auxiliary space.
\end{enumerate}
\section{Example: binary search}
Assume $A[0\ldots n-1]$ is sorted in nondecreasing order. Maintain half-open $[l,r)$ containing every index that could hold target $x$. Initially $l=0,r=n$. Choose $m=l+\lfloor (r-l)/2\rfloor$. If $A[m]<x$ set $l=m+1$; otherwise set $r=m$. The interval narrows; at termination $l=r$ is the first position not less than $x$. Check equality before reporting found. The variant $r-l$ decreases and is never negative.
\begin{engineeringnote}[title=Implementation warning]
The code computes the midpoint from the interval width, rather than adding two potentially large nonnegative indices before division. Language arithmetic still matters: C++ size types and signed conversions require attention.
\end{engineeringnote}
\section{Example: amortized vector growth}
Suppose capacity doubles when full. A sequence of $n$ appends copies at most $1+2+4+\cdots+2^{\lfloor\log_2 n\rfloor}<2n$ existing entries; adding the $n$ writes gives fewer than $3n$ primitive element operations (up to convention). Thus amortized append cost is $O(1)$, but any resize may be $\Theta(n)$.
\section{Probabilistic and experimental statements}
\begin{itemize}\item Worst case describes every admissible input; expected describes an explicit probability model.
\item A p-value or test failure cannot certify true randomness, and sampling cannot establish cryptographic unpredictability.
\item A measured trend is not a complexity proof. Fit candidate curves only as diagnostics after deriving theoretical expectations.
\item A confidence interval for mean latency is not automatically a confidence interval for an individual future execution.
\end{itemize}
\section{Core recurrence laboratory}
For $T(n)=T(n-1)+n$, telescope to $\sum_{k=1}^{n}k=\Theta(n^2)$. For $T(n)=2T(n/2)+n$, the level argument gives $\Theta(n\log n)$. For $T(n)=T(n/2)+1$, there are $\Theta(\log n)$ levels. In every proof specify $T(1)$ and whether $n$ is restricted to powers of two.
\section{Independent exercises}
\begin{enumerate}
\item Show an invariant for a stable merge operation and explain when two equal-key elements can change relative order.
\item Derive expected comparisons for a linear scan under a uniformly selected successful key position; distinguish unsuccessful queries.
\item Prove that recursive enumeration of all subsets has at least $2^n$ leaf visits, then explain why the recursion stack can be only $O(n)$.
\item Construct a comparator that violates transitivity and explain why a sorting algorithm's contract fails for it.
\item Write a performance claim using separate asymptotic, allocation and cache-behavior evidence.
\end{enumerate}
\SourceNote{Vol.~1, \S\S1.2.1, 1.2.10--1.2.11; Vol.~3, \S\S5.2, 6.2.1; Vol.~4A, \S7.2.1.}
\end{document}
